\documentclass[11pt]{article}
\usepackage[utf8]{inputenc}
\usepackage[top=0.6in,bottom=0.6in,left=0.7in,right=0.7in]{geometry}
\usepackage{graphicx}
\usepackage{booktabs}
\usepackage{array}
\usepackage{hyperref}
\usepackage{amsmath}
\usepackage{caption}
\usepackage{float}
\usepackage{xcolor}
\usepackage{multirow}
\usepackage{fancyhdr}
\usepackage{titling}

\pagestyle{fancy}
\fancyhf{}
\fancyhead[L]{\small Moshiko 2.0 Project Report}
\fancyhead[R]{\small Group: Member 1, Member 2, Member 3}
\fancyfoot[C]{\thepage}
\renewcommand{\headrulewidth}{0.8pt}
\renewcommand{\footrulewidth}{0.4pt}

\setlength{\parindent}{0pt}
\setlength{\parskip}{4pt}
\setlength{\abovecaptionskip}{4pt}
\setlength{\belowcaptionskip}{2pt}

\title{\Large\textbf{Moshiko 2.0: Quantization Analysis and Speech-to-Speech Model Optimization}}
\author{\small
    Member 1 (ID: XXXXXXXXX) \quad
    Member 2 (ID: XXXXXXXXX) \quad
    Member 3 (ID: XXXXXXXXX)
}
\date{\small\today}

\begin{document}

\maketitle
\vspace{-0.3in}

\section*{Abstract}
We present a systematic quantization study of the Moshiko 7B speech-text model (Kyutai) on a Google Colab Tesla T4 GPU (16 GB VRAM). Our experiments show Kyutai's INT8 quantization is lossless (0.00 dB SNR difference from BF16), the Mimi codec operates 30--140$\times$ faster than real-time (RTF: 0.007--0.031), and decoder convolutional layers are the most quantization-sensitive. We also implemented a mixed-precision strategy retaining the top 20\% sensitive layers at BF16.

\section{Project Overview}
The Moshiko project evaluates the Moshiko speech-text model by Kyutai (\url{https://github.com/kyutai-labs/moshi}), a 7B-parameter multimodal model comprising:

\begin{itemize}[noitemsep,topsep=2pt]
    \item \textbf{Mimi Audio Codec:} Neural codec encoding 24 kHz mono audio to discrete tokens $[B, 8, T]$ and decoding them back.
    \item \textbf{Moshiko LM:} 7B transformer (32 layers, 4096 dim, 32 heads) processing speech and text simultaneously.
\end{itemize}

\textbf{Original Plan:} We initially aimed to run \textbf{two full-duplex speech models simultaneously} (Moshi + another model). However, the BF16 LM alone requires 15.4 GB VRAM, the architecture uses custom \texttt{weights\_per\_step} streaming, and dual-model documentation was unavailable. We pivoted to quantization analysis and fine-tuning optimization.

\textbf{Revised Plan:} (1) \textbf{Quantization (First Demo):} Compare BF16/Q8, apply INT8/INT4 PTQ, layer-wise sensitivity, mixed-precision. (2) \textbf{Fine-Tuning (Future):} LoRA (train 0.1\% params) or Frozen Backbone + New Head (e.g., language ID classifier).

\textbf{What We Did:} Set up PyTorch 2.4.0 + CUDA 12.1, loaded Mimi (q8) + Moshiko LM (q8) within 8.4--8.8 GB VRAM, ran codec tests, latency benchmarks, BF16 vs Q8 comparison, and layer-wise sensitivity analysis.

\textbf{What's Next:} Complete speech-to-speech generation (Test 3), finish LM INT4 quantization, implement language ID fine-tuning, and explore GPTQ/AWQ.

\section{Technical Setup}
Platform: Google Colab (free tier), NVIDIA Tesla T4 (15.6 GB VRAM, CUDA 12.1), Python 3.12.13, PyTorch 2.4.0, \texttt{moshi} v0.2.13, bitsandbytes v0.49.2, safetensors v0.7.0. Model files stored on Google Drive.

\section{Results}

\subsection{Mimi Codec Encode/Decode Quality}
\begin{table}[H]
\centering
\caption{Codec performance on synthetic and real audio.}
\label{tab:codec_quality}
\begin{tabular}{@{}lccc@{}}
\toprule
\textbf{Test Signal} & \textbf{Encode (ms)} & \textbf{Decode (ms)} & \textbf{SNR (dB)} \\
\midrule
440 Hz Sine Wave (3s) & 9,618.8 & 119.7 & 15.88 \\
Real Speech (12s WAV) & 2,186.6 & 25.3 & 6.20 \\
\bottomrule
\end{tabular}
\end{table}

\subsection{Latency Benchmark}
\begin{table}[H]
\centering
\caption{Latency and real-time factor (RTF) across durations.}
\label{tab:latency}
\begin{tabular}{@{}lcccc@{}}
\toprule
\textbf{Duration} & \textbf{Encode (ms)} & \textbf{Decode (ms)} & \textbf{Total (ms)} & \textbf{RTF} \\
\midrule
1s & 16.3 & 14.9 & 31.2 & 0.031 \\
3s & 18.4 & 17.2 & 35.5 & 0.012 \\
5s & 17.2 & 21.0 & 38.2 & 0.008 \\
10s & 35.0 & 34.7 & 69.7 & 0.007 \\
\bottomrule
\end{tabular}
\end{table}

All RTF values are well below 1.0---the codec operates \textbf{significantly faster than real-time}.

\subsection{BF16 vs. Q8 Comparison}
\begin{table}[H]
\centering
\caption{BF16 vs. INT8 (q8) Mimi codec SNR comparison.}
\label{tab:bf16_vs_q8}
\begin{tabular}{@{}lccc@{}}
\toprule
\textbf{Signal Type} & \textbf{Q8 SNR (dB)} & \textbf{BF16 SNR (dB)} & \textbf{Difference (dB)} \\
\midrule
Pure Tone 440 Hz & 15.43 & 15.43 & 0.00 \\
Mixed Tones & 7.44 & 7.44 & 0.00 \\
White Noise & -1.54 & -1.54 & 0.00 \\
Chirp (Sweep) & 0.68 & 0.68 & 0.00 \\
\bottomrule
\end{tabular}
\end{table}

\textbf{Key Finding:} Kyutai's INT8 quantization is \textbf{completely lossless}---bit-identical output across all signals.

\subsection{Layer-Wise Sensitivity Analysis}
We quantized each of 130 Mimi codec layers individually to INT8. Total quantizable parameters: 79,292,225. BF16 baseline SNR: 15.88 dB.

\begin{table}[H]
\centering
\caption{Top 5 most quantization-sensitive layers.}
\label{tab:sensitive_layers}
\begin{tabular}{@{}clcc@{}}
\toprule
\textbf{Rank} & \textbf{Layer} & \textbf{Type} & \textbf{SNR Drop (dB)} \\
\midrule
1 & \texttt{decoder.model.0.conv.conv} & Conv1d & 2.33 \\
2 & \texttt{encoder.model.14.conv.conv} & Conv1d & 0.59 \\
3 & \texttt{encoder.model.4.block.3.conv.conv} & Conv1d & 0.47 \\
4 & \texttt{encoder.model.9.conv.conv} & Conv1d & 0.47 \\
5 & \texttt{decoder.model.3.block.1.conv.conv} & Conv1d & 0.26 \\
\bottomrule
\end{tabular}
\end{table}

\subsection{Mixed-Precision Strategy}
Top 20\% sensitive layers (19) kept at BF16; remaining 80\% (79) quantized to INT8.

\begin{table}[H]
\centering
\caption{Mixed-precision vs. baselines.}
\label{tab:mixed_precision}
\begin{tabular}{@{}lcc@{}}
\toprule
\textbf{Model Variant} & \textbf{SNR (dB)} & \textbf{Drop from Baseline (dB)} \\
\midrule
BF16 (Baseline) & 15.88 & 0.00 \\
Kyutai Q8 & 15.88 & 0.00 \\
Mixed (80\% INT8) & 17.31 & -1.43 \\
\bottomrule
\end{tabular}
\end{table}

Mixed model file size: 384.8 MB. The higher SNR indicates different (not necessarily worse) reconstruction characteristics.

\subsection{INT4 Quantization (NF4)}
Applied bitsandbytes NF4 quantization. Conv1d layers could not be quantized (API limitation: \texttt{quantize\_4bit} returned 2 values instead of 3). INT4 model saved as safetensors (storage-only; requires dequantization before inference).

\subsection{VRAM Usage}
\begin{table}[H]
\centering
\caption{GPU memory consumption by component.}
\label{tab:vram}
\begin{tabular}{@{}lcl@{}}
\toprule
\textbf{Component} & \textbf{VRAM} & \textbf{Notes} \\
\midrule
Mimi Codec (q8) & $\sim$0.2 GB & Audio codec \\
Moshiko LM (q8) & $\sim$8.4 GB & 7B model, INT8 \\
Moshiko LM (bf16) & $\sim$15.4 GB & Does not fit on T4 \\
Activation (inference) & $\sim$2--4 GB & Depends on input length \\
\midrule
\textbf{Total (q8)} & \textbf{$\sim$12--14 GB} & \textbf{Fits on T4} \\
\textbf{Total (bf16)} & \textbf{$\sim$18+ GB} & \textbf{Does not fit on T4} \\
\bottomrule
\end{tabular}
\end{table}

\subsection{PyTorch Dynamic Quantization Incompatibility}
\texttt{torch.quantization.quantize\_dynamic()} is \textbf{incompatible} with moshi's custom transformer architecture (\texttt{weights\_per\_step} scheduling). Manual INT8 weight quantization (round-to-nearest, per-tensor scaling) was implemented instead.

\section{Quantization Experiment Status}
\begin{table}[H]
\centering
\caption{Summary of all quantization experiments.}
\label{tab:exp_status}
\begin{tabular}{@{}llc@{}}
\toprule
\textbf{Experiment} & \textbf{Type} & \textbf{Status} \\
\midrule
1: BF16 vs. Q8 & Comparison & Complete \\
2a: Manual INT8 PTQ & Post-Training Quantization & Partial \\
2b: INT4 PTQ (NF4) & Post-Training Quantization & Complete \\
3a: Layer-wise Sensitivity & Sensitivity Analysis & Complete \\
3b: Mixed Precision & Strategy & Complete \\
4: Dynamic Quantization & Runtime Quantization & Planned \\
5a: LM Layer Sensitivity & Sensitivity Analysis & Started \\
5b: LM INT4 (NF4) & Post-Training Quantization & Planned \\
6: Final Report & Summary & Planned \\
\bottomrule
\end{tabular}
\end{table}

\section{Discussion and Key Findings}
\begin{enumerate}[noitemsep,topsep=2pt]
    \item \textbf{Kyutai's Q8 is lossless:} 0.00 dB SNR difference from BF16 across all signals---exceptionally well-designed quantization.
    \item \textbf{Mimi codec is highly efficient:} RTF as low as 0.007 (140$\times$ faster than real-time).
    \item \textbf{Decoder conv layers are most sensitive:} First decoder conv layer drops 2.33 dB when quantized---critical for mixed-precision.
    \item \textbf{Standard quantization APIs fail:} PyTorch dynamic quantization incompatible; manual implementation required.
    \item \textbf{T4 is sufficient for q8:} 8.4--8.8 GB VRAM usage enables 7B model on free-tier Colab.
\end{enumerate}

\textbf{Limitations:} BF16 LM unavailable on T4 (15.4 GB $>$ 15.6 GB); speech-to-speech generation (Test 3) hung during LM generation; INT4 inference requires dequantization step; bitsandbytes does not support Conv1d NF4.

\section{Future Work}
\begin{itemize}[noitemsep,topsep=2pt]
    \item \textbf{Complete Test 3:} Use official \texttt{InferenceState} API (Option C) or LMGen float32 fix (Option B).
    \item \textbf{Complete quantization:} Execute Experiments 4, 5b, 6 (dynamic quantization, LM INT4, final dashboard).
    \item \textbf{Language ID fine-tuning:} Add frozen-LM classification head (Linear(4096, 512) $\rightarrow$ ReLU $\rightarrow$ Dropout(0.3) $\rightarrow$ Linear(512, $N$)). Estimated 2--4 hours on T4.
    \item \textbf{Dialogue fine-tuning:} Train on 100+ hours paired speech data. AdamW, lr $\sim 10^{-5}$, 3--5 days/epoch on T4.
    \item \textbf{LoRA:} Train 0.1\% of parameters for task-specific adaptation.
    \item \textbf{Advanced quantization:} GPTQ, AWQ, INT2, fine-tuning quantized models.
\end{itemize}

\section{Conclusion}
We successfully set up and evaluated the Moshiko 7B model on a T4 GPU through systematic quantization. Key contributions: (1) demonstrated Kyutai's INT8 quantization is lossless, (2) benchmarked Mimi codec at RTF 0.007--0.031, (3) identified decoder conv layers as most quantization-sensitive, (4) implemented a mixed-precision strategy, and (5) documented PyTorch quantization API incompatibility. The project evolved from a dual-model plan to a focused quantization study, providing a foundation for fine-tuning and advanced compression work.

\begin{thebibliography}{9}
\bibitem{moshi} Kyutai Labs. Moshi: A Speech-Text Foundation Model. \url{https://github.com/kyutai-labs/moshi}
\bibitem{moshiko} Kyutai Labs. Moshiko PyTorch Models. \url{https://huggingface.co/kyutai/moshiko-pytorch-bf16}
\bibitem{bnb} Dettmers, T. et al. bitsandbytes. \url{https://github.com/TimDettmers/bitsandbytes}
\bibitem{lora} Hu, E.J. et al. LoRA. \textit{arXiv:2106.09685}, 2021.
\bibitem{gptq} Frantar, E. et al. GPTQ. \textit{arXiv:2210.17323}, 2022.
\bibitem{awq} Lin, J. et al. AWQ. \textit{arXiv:2306.00978}, 2023.
\end{thebibliography}

\end{document}
